% HandoutsTexTemplate
% Copyright (C) 2026 Giulio Salvi
%
% This is free software: you can redistribute it and/or modify
% it under the terms of the GNU General Public License as published by
% the Free Software Foundation, either version 3 of the License, or
% (at your option) any later version.

% This software is distributed in the hope that it will be useful,
% but WITHOUT ANY WARRANTY; without even the implied warranty of
% MERCHANTABILITY or FITNESS FOR A PARTICULAR PURPOSE.  See the
% GNU General Public License for more details.
%
% You should have received a copy of the GNU General Public License
% along with this program.  If not, see <https://www.gnu.org/licenses/>.

% Match hyperlink colors to the template accent palette.
\hypersetup{
  colorlinks=true,
  linkcolor=linkcolor,
  urlcolor=linkcolor,
  citecolor=linkcolor
}

% Cover font uses normal, readable interword spacing.
\NewDocumentCommand{\covertitle}{m}{%
  {\Huge\bfseries #1\par}%
}

% Keep box titles and counter tags on the same line.
\newcommand{\boxtitle}[2]{#1\hspace{0.55em}#2}
\newcommand{\boxcountertag}[1]{%
  \hfill\nobreak\mbox{\normalfont #1}%
}
\newcommand{\boxtitlewithlabel}[3]{%
  \boxtitle{#1}{#2}\boxcountertag{#3}%
}
\newcommand{\boxplainwithlabel}[2]{%
  #1\boxcountertag{#2}%
}
\newcommand{\boxtitlewithname}[4]{%
  \ifstrempty{#3}{%
    \boxtitlewithlabel{#1}{#2}{#4}%
  }{%
    \boxtitlewithlabel{#1}{#2 of #3}{#4}%
  }%
}
\newcommand{\boxtitleorfallback}[4]{%
  \ifstrempty{#3}{%
    \boxtitlewithlabel{#1}{#2}{#4}%
  }{%
    \boxtitlewithlabel{#1}{#3}{#4}%
  }%
}

% A deliberate blank page is used only behind the cover.
\NewDocumentCommand{\blankpage}{}{%
  \thispagestyle{empty}\null\clearpage
}

% Native book opening policy, also used by dedication and epigraph pages.
\makeatletter
\NewDocumentCommand{\handoutsclearchapterpage}{}{%
  \if@openright\cleardoublepage\else\clearpage\fi
}
\makeatother

% Visible labels follow Babel's selected main language.
\NewDocumentCommand{\defaultcolophontext}{}{These notes were typeset using \LaTeX{}.}
\NewDocumentCommand{\prefacename}{}{Preface}
\NewDocumentCommand{\licensename}{}{License}
\addto\captionsitalian{%
  \renewcommand{\defaultcolophontext}{Queste dispense sono state composte con \LaTeX{}.}%
  \renewcommand{\prefacename}{Prefazione}%
  \renewcommand{\licensename}{Licenza}%
}
\addto\captionsenglish{%
  \renewcommand{\defaultcolophontext}{These notes were typeset using \LaTeX{}.}%
  \renewcommand{\prefacename}{Preface}%
  \renewcommand{\licensename}{License}%
}

% Draw a simple rectangular border on the cover page.
\newcommand{\drawcoverborder}{%
  \begin{tikzpicture}[remember picture,overlay]
    \draw[
      black!90,
      line width=1pt
    ]
      ($(current page.north west)+(1.8cm,-1.8cm)$)
      rectangle
      ($(current page.south east)+(-1.8cm,1.8cm)$);
  \end{tikzpicture}%
}

% Real chapters retain the original large opening titles.
\NewDocumentCommand{\headertitle}{m}{\textit{#1}}
\titleformat{\chapter}[display]
  {\normalfont\raggedright}
  {\bfseries\LARGE\chaptertitlename~\thechapter}
  {2.5ex}
  {\bfseries\Huge}
\titlespacing*{\chapter}{0pt}{0.10\textheight}{0.06\textheight}
\titleformat{\section}{\Large\bfseries}{\thesection}{0.7em}{}
\titlespacing*{\section}{0pt}{2.5ex plus 1ex minus 0.2ex}{1ex plus 0.2ex}
\titleformat{\subsection}{\large\bfseries}{\thesubsection}{0.7em}{}
\titlespacing*{\subsection}{0pt}{2ex plus 0.8ex minus 0.2ex}{0.8ex plus 0.2ex}

% Chapters and sections feed the two running headers; starred headings stay native.
\makeatletter
\renewcommand{\chaptermark}[1]{%
  \markboth{\headertitle{%
    \ifnum\value{secnumdepth}>-1\relax
      \if@mainmatter\chaptername~\thechapter.\ \fi
    \fi
    #1%
  }}{}%
}
\makeatother
\renewcommand{\sectionmark}[1]{%
  \markright{\headertitle{%
    \ifnum\value{secnumdepth}>0\relax\thesection.\ \fi
    #1%
  }}%
}

% Opening pages have a centred footer; blank versos are handled by emptypage.
\fancypagestyle{plain}{%
  \fancyhf{}
  \fancyfoot[C]{\thepage}
  \renewcommand{\headrulewidth}{0pt}
  \renewcommand{\footrulewidth}{0pt}
}
\fancypagestyle{mainmatter}{%
  \fancyhf{}
  \fancyhead[LE]{\small\thepage}
  \fancyhead[RE]{\parbox[b]{0.82\textwidth}{\raggedleft\small\nouppercase{\leftmark}}}
  \fancyhead[LO]{\parbox[b]{0.82\textwidth}{\raggedright\small\nouppercase{\rightmark}}}
  \fancyhead[RO]{\small\thepage}
  \renewcommand{\headrulewidth}{0pt}
  \renewcommand{\footrulewidth}{0pt}
}
\fancypagestyle{frontmattersection}{%
  \fancyhf{}
  \fancyhead[LE]{\small\thepage}
  \fancyhead[RE]{\parbox[b]{0.82\textwidth}{\raggedleft\small\nouppercase{\leftmark}}}
  \fancyhead[LO]{\parbox[b]{0.82\textwidth}{\raggedright\small\nouppercase{\leftmark}}}
  \fancyhead[RO]{\small\thepage}
  \renewcommand{\headrulewidth}{0pt}
  \renewcommand{\footrulewidth}{0pt}
}

% Forward optional titles and options to box environments.
\NewDocumentCommand{\beginbox}{m O{} g}{%
  \ifstrempty{#2}{%
    \IfNoValueTF{#3}{\begin{#1}}{\begin{#1}{#3}}%
  }{%
    \IfNoValueTF{#3}{\begin{#1}[#2]}{\begin{#1}[#2]{#3}}%
  }%
}
\NewDocumentCommand{\endbox}{m}{\end{#1}}

% Mathtools paired delimiters keep the existing auto-sized shorthand syntax.
\NewDocumentCommand{\bal}{m}{\mathrm{bal}(#1)}
\DeclarePairedDelimiter{\handoutsset}{\lbrace}{\rbrace}
\DeclarePairedDelimiter{\handoutsabs}{\lvert}{\rvert}
\DeclarePairedDelimiter{\handoutsnorm}{\lVert}{\rVert}
\NewDocumentCommand{\set}{m}{\handoutsset*{#1}}
\NewDocumentCommand{\abs}{m}{\handoutsabs*{#1}}
\NewDocumentCommand{\norm}{m}{\handoutsnorm*{#1}}

% Build the shared badge style used by custom lists.
\newcounter{numberedlistitemctr}
\newlength{\numberedlistitembadgewidth}
\newcommand{\numberedlistitembadge}[1]{%
  \tikz[baseline=(n.base)]{
    \node[
      fill=linkcolor,
      text=white,
      font=\bfseries\scriptsize,
      minimum width=1em,
      minimum height=1em,
      inner sep=0pt,
      rounded corners=1pt
    ] (n) {#1};
  }%
}
\newcommand{\unorderedlistitembadge}{%
  \numberedlistitembadge{\textbullet}%
}

% Keep numbered and unordered badge lists aligned the same way.
\newcommand{\badgelistlayout}{%
  \setlength{\labelwidth}{\numberedlistitembadgewidth}
  \setlength{\labelsep}{0.75em}
  \setlength{\leftmargin}{\dimexpr\labelwidth+\labelsep\relax}
  \setlength{\itemsep}{0.45em}
  \setlength{\topsep}{0pt}
  \setlength{\partopsep}{0pt}
  \setlength{\parsep}{0pt}
  \setlength{\listparindent}{0pt}
  \setlength{\itemindent}{0pt}
  \setlength{\rightmargin}{0pt}
}
\settowidth{\numberedlistitembadgewidth}{\numberedlistitembadge{99}}

% Numbered list with boxed numeric badges.
\NewDocumentEnvironment{numberedlist}{}{%
  \setcounter{numberedlistitemctr}{0}%
  \begin{list}{\numberedlistitembadge{\arabic{numberedlistitemctr}}}{%
    \usecounter{numberedlistitemctr}
    \badgelistlayout
  }%
  \ignorespaces
}{%
  \end{list}\ignorespacesafterend
}

% Unordered list with the same visual badge style.
\NewDocumentEnvironment{unorderedlist}{}{%
  \begin{list}{\unorderedlistitembadge}{%
    \badgelistlayout
  }%
  \ignorespaces
}{%
  \end{list}\ignorespacesafterend
}

% Derivations are just titled numbered badge lists.
\NewDocumentEnvironment{derivation}{O{\derivationname:}}{%
  \par\medskip
  \noindent\textbf{#1}\par\smallskip
  \begin{numberedlist}
}{%
  \end{numberedlist}\par\medskip
}
\newcommand{\derivstep}[2]{%
  \item $\displaystyle #1$\hfill\nobreak\mbox{\textit{(#2)}}%
}

% Proof cases align a short label on the left and indent the full body.
\newlength{\proofcaselabelwidth}
\settowidth{\proofcaselabelwidth}{$(\Leftarrow)$}
\newcommand{\proofcaselabel}[1]{%
  \makebox[\proofcaselabelwidth][l]{$#1$}%
}
\NewDocumentEnvironment{proofcase}{m}{%
  \par\medskip
  \begin{list}{}{%
    \setlength{\labelwidth}{\proofcaselabelwidth}
    \setlength{\labelsep}{0.45em}
    \setlength{\leftmargin}{\dimexpr\labelwidth+\labelsep\relax}
    \setlength{\itemsep}{0pt}
    \setlength{\topsep}{0pt}
    \setlength{\partopsep}{0pt}
    \setlength{\parsep}{0pt}
    \setlength{\listparindent}{0pt}
    \setlength{\itemindent}{0pt}
    \setlength{\rightmargin}{0pt}
  }%
  \item[\proofcaselabel{#1}]%
  \ignorespaces
}{%
  \end{list}\ignorespacesafterend
}

% Wrap a centered figure with a caption.
\newcommand{\img}[3][]{%
  \begin{figure}[H]
    \centering
    \includegraphics[#1]{#2}
    \caption{#3}
  \end{figure}
}
